\documentclass[12pt]{article}
\usepackage[utf8]{inputenc}
\usepackage[T1]{fontenc}
\usepackage{amsmath, amssymb}
\usepackage{geometry}
\geometry{margin=2.5cm}
\setlength{\emergencystretch}{3em}
\renewcommand{\abstractname}{Özet}
\title{\textbf{Duran Sicimler Teoremi: Zamanın ve Entropinin Yokluğunda Mükemmel Düzenin Yeniden Doğuşu}}
\author{Fahri Meral — 2025}
\date{}

\begin{document}
\maketitle

\begin{abstract}
Bu çalışma, sicimlerin titreşimi ile zaman, entropi ve algı arasındaki ilişkiyi yeniden yorumlamaktadır. 
Klasik sicim kuramı, evrenin temel yapıtaşlarının titreşen sicimler olduğunu savunur. 
Ancak bu teorem, insanların yalnızca titreşen sicimleri algılayabildiğini ve titreşimin durması durumunda ortaya çıkan 
"mükemmel düzeni" fark edemediğini öne sürmektedir. 
Bu bakış açısına göre kara delikler, sicimlerin titreşimini durdurarak evreni yeniden mükemmelliğe ulaştıran geçitlerdir. 
Dolayısıyla kara delikler, yok oluş değil, farklı bir boyuta geçişin fiziksel yansımasıdır.
\end{abstract}

\section{Giriş}
Evrenin kökenine dair modern fizik, Büyük Patlama (Big Bang) ile başlayan bir zaman oku varsayar. 
Bu an, sicimlerin ilk titreşimlerinin başladığı, yani enerjinin düzenli hâlden düzensiz hâle geçtiği noktadır. 
Bu süreçte entropi artışı, mükemmel olandan sapmanın kozmik ölçekteki ifadesidir.

\section{Duran Sicim Varsayımı}
Bu teorem, sicimlerin titreşiminin durduğu durumda zamanın ve entropinin anlamını yitirdiğini öne sürer. 
Titreşimin sıfırlandığı (\(\omega = 0\)) durumda sicimler mükemmel bir düzen oluşturur. 
Bu durumda enerji değişimi, dalga yayılımı ve hatta foton hareketi ortadan kalkar. 
İnsanın algı kapasitesi, yalnızca titreşen sistemleri (enerjiyi) gözlemleyebildiğinden, bu mükemmel düzen “karanlık” olarak görünür.

\section{Kara Delikler ve Boyutsal Geçiş}
Kara delikler genellikle yoğun kütle çekimi nedeniyle ışığın bile kaçamadığı bölgeler olarak tanımlanır. 
Ancak bu teorem, kara deliklerin özünde sicimlerin titreşimini durduran yapılar olduğunu öne sürer. 
Bu durumda kara delikler, enerji ve zamanın bittiği, mükemmelliğin yeniden sağlandığı geçitlerdir. 
Evrenin her kara deliği, başka bir mükemmel düzen boyutuna açılan bir “kapı” olarak işlev görür.

\section{Kozmolojik Sonuçlar}
Bu yaklaşım, evrenin başlangıcında saf ve mükemmel bir düzenin bulunduğunu, 
zamanın ve entropinin sicimlerin titreşmesiyle ortaya çıktığını ve 
kara deliklerin bu süreci tersine çevirdiğini savunur. 
Dolayısıyla evrenin “sonu”, bir çöküş değil, 
başlangıçtaki mükemmel hâlin yeniden kazanılması anlamına gelir.

\section{Sonuç}
Duran Sicimler Teoremi, fiziksel evrenin algılanabilir sınırlarının ötesinde, 
titreşimin durduğu bir mükemmel düzenin var olabileceğini öne sürmektedir. 
Bu düzen, ne enerjiye ne de zamana bağlıdır. 
İnsanın algısı bu boyuta erişemediğinden, kara delikler bize karanlık görünür; 
ancak özünde, onlar evrenin yeniden doğduğu alanlardır.
\end{document}
